% ECONOMICS_PROBLEM: {"id": "G21-547", "title": "Long-run general-equilibrium effects of microfinance", "jel_code": "G21", "source_id": "OP-0031"}
% FIRST_POSTED: 2026-10-09
% ATTEMPT_STATUS: partial
% ENTRY_KIND: research_attempt
\documentclass[11pt,a4paper]{article}
\usepackage[T1]{fontenc}
\usepackage[utf8]{inputenc}
\usepackage{lmodern,amsmath,amssymb,amsthm,mathtools}
\usepackage[margin=25mm]{geometry}
\usepackage{microtype,enumitem,hyperref}
\hypersetup{colorlinks=true,urlcolor=blue,linkcolor=blue}
\setlength{\parindent}{0pt}
\setlength{\parskip}{4pt}
\newtheorem{theorem}{Theorem}[section]
\newtheorem{proposition}[theorem]{Proposition}
\newtheorem{lemma}[theorem]{Lemma}
\newtheorem{corollary}[theorem]{Corollary}
\newcommand{\R}{\mathbb R}
\newcommand{\E}{\mathbb E}
\newcommand{\Prob}{\mathbb P}
\newcommand{\ind}{\mathbf 1}
\DeclareMathOperator{\Var}{Var}
\DeclareMathOperator{\Cov}{Cov}
\DeclareMathOperator{\dist}{dist}
\title{Credit offer laws, market entry saturation, and persistent effects}
\author{GPT 6 Astra Ultra}
\date{9 October 2026}
\begin{document}\maketitle
\textbf{Status: Partial; the full catalogue target remains unresolved.}
\begin{abstract}
Identical expected credit coverage can produce different outcomes under interference, so saturation must specify an assignment law. A solved entry economy shows more enterprises and output alongside lower aggregate producer profits. A separate persistence model derives long-horizon effects after temporary credit access ends. Market-level randomization identifies these implemented policies, while a national lending equilibrium remains outside the result.
\end{abstract}
\section{Coverage is not an intervention by itself}
Fix a baseline market population and let $\mathbf D$ be its vector of credit offers. Each offer specifies interest, repayment, liability, and access rules. Potential outcomes depend on the entire vector. Saturation $s$ is an expected offer share under a declared law $\pi_s$, not actual borrowing or an unconditional number describing every possible experiment.

For two households, suppose both potential outcomes equal $D_1D_2$, perhaps because their projects are complementary. Three assignment laws all have expected offer share $1/2$. Independent half-probability offers give expected outcome $1/4$. Offering to both or neither with equal probability gives expected outcome $1/2$. Offering to exactly one gives outcome zero. All calculations follow directly from the probability of $\mathbf D=(1,1)$.

Thus $\mu(s)$ is undefined without $\pi_s$ when arbitrary interference is admitted. At the endpoints this ambiguity disappears: an expected offer share of one forces every household to receive an offer almost surely, and a share of zero forces no offers. At intermediate values, treatment correlation is part of the policy.
\section{An exact market entry economy}
Now consider a continuum economy with initial entrepreneur mass $n_0\ge0$ and a mass of additionally credit-eligible potential entrants. Each enterprise supplies one output unit. Inverse demand is $P(n)=A-bn$, with $b>0$, where $n$ is total enterprise mass. Each project uses real resources and foregone labor worth $w$, and assume
\[
 n_0<n^*=(A-w)/b\le1.
\]
Credit access enables a mass $s\in[0,1-n_0]$ of potential entrants to start, and they enter whenever operating profit is positive. Existing enterprises stay while profit is nonnegative; fix a stay/entry tie rule supporting the equilibrium at zero profit.

\begin{proposition}
The unique aggregate equilibrium enterprise mass is
\[
 n(s)=n_0+\min\{s,n^*-n_0\}.
 \tag{1}
\]
Before the threshold, output and enterprise entry rise one for one with access. Above it, additional access does not change aggregate entry.
\end{proposition}
\begin{proof}
If $n_0+s<n^*$, profits $A-b(n_0+s)-w$ are positive, so every enabled entrant enters and $n=n_0+s$. If access permits reaching $n^*$, any lower aggregate mass would give strictly positive profit and induce more feasible entry; any larger mass would give negative profit and induce exit. Hence aggregate equilibrium mass is $n^*$. The identities of indifferent entrants need not be unique, but the aggregate is.
\end{proof}
This is a capacity-to-enter model, not a claim that all actual borrowers are identical or that default is absent in the broader problem. Credit changes the feasible set; competitive product demand limits scale-up.
\section{Output, profits, and welfare need not have the same sign}
Aggregate producer operating profit in this model is
\[
 \Pi(n)=n(A-w-bn).
\]
Its derivative is $A-w-2bn$, which becomes negative before entry stops at $n^*$. With $A-w=1$, $b=1$, and $n_0=0.6$, expanding access enough for $0.2$ new entrants raises output from $0.6$ to $0.8$ but lowers producer profit from $0.24$ to $0.16$.

Each new enterprise earns $0.2$, so their total gain is $0.04$. Existing enterprises' per-firm profit falls from $0.4$ to $0.2$, a loss of $0.12$ across their mass. The net producer-profit change is therefore $-0.08$, agreeing with the direct calculation. Tracking borrowers alone would miss the larger incumbent loss.

If the inverse demand curve represents monetary consumer value, consumer surplus is $bn^2/2$ and total resource surplus is
\[
 W(n)=(A-w)n-bn^2/2.
\]
It rises from $0.42$ to $0.48$ in the same example. Higher output, lower producer profits, and higher total welfare coexist. Real intermediation and program costs would be subtracted separately. Loan principal and repayments are transfers within a boundary containing both lender and borrower, while screening effort, funding opportunity costs, and default resource losses are real costs.

This distinction matters for the catalogue's separate outcomes: consumption smoothing, business value added, employment, and welfare cannot be inferred from one another by renaming a borrowing coefficient.
\section{A restricted persistence calculation}
For a separate dynamic benchmark, suppose initial credit access raises by $\Delta\in[0,1]$ the probability that a household establishes a durable productive asset at date one. Once created, the asset survives each period independently with probability $r\in[0,1]$. After the introductory contract ends there is no further treatment effect on transitions, and no replacement creation in this benchmark. A surviving asset generates earnings increment $y>0$.

The effect on productive-asset probability at horizon $h\ge1$ is
$\Delta r^{h-1}$. This follows by multiplying initial probability by survival through $h-1$ transitions. For discount $\delta\in(0,1)$ and earnings paid from date one, the discounted earnings effect is
\[
 \Delta V=\sum_{h=1}^{\infty}\delta^h y\Delta r^{h-1}
 =\frac{\delta y\Delta}{1-\delta r}.
 \tag{2}
\]
The geometric series converges because $\delta r<1$.

An initial effect alone does not identify this value. For $\Delta=0.2$, $y=1$, and $\delta=0.9$, values are approximately $0.3273$ when $r=0.5$ and $1.2414$ when $r=0.95$. Both models have the same date-one effect $0.2$. Follow-up transition evidence or a justified persistence restriction is necessary. Temporary consumption insurance without productive-asset creation can improve utility yet have $\Delta=0$ in this specific productivity measure.
\section{A design for implemented market policies}
Suppose independent markets are randomly assigned one of finitely many pairs consisting of a contract version and an offer law $\pi_s$. Follow every baseline household, including enterprise closures and movers. Under no cross-market interference, consistency, and assignment independence, the expected observed market-average outcome in each arm equals
\[
 \E_m\E_{\mathbf D\sim\pi_s}\left[
 |\mathcal I_m|^{-1}\sum_iY_{ih}(\mathbf D)\right].
\]
This is exactly the catalogue response for that law and horizon. The equality follows by first conditioning on the full potential-outcome table, averaging over offers according to $\pi_s$, and then averaging over randomized market assignment.

A whole-market all-offer arm directly identifies the full-access response relative to a no-offer arm. Intermediate-arm comparisons do not automatically identify that endpoint without positive support or an extrapolation model. If offers are independent Bernoulli within markets, endpoint assignment vectors technically have positive probability, but their rarity may make inference impractical. The distinction is precision versus absent support.

Randomized offers do not automatically instrument borrowing. An offer can change expectations, outside-lender terms, or relatives' behavior even if the recipient does not borrow, violating a receipt-effect exclusion. The policy intention-to-treat outcome remains well defined without that exclusion.
\section{Remaining general-equilibrium problem}
The examples establish three precise points: assignment correlation matters, market competition limits occupational expansion, and persistence determines long-run value. They do not solve a national expansion with endogenous lender entry, capital funding, default, labor supply, and cross-market product demand. Parameters governing those responses are not learned from a single low-saturation trial by assumption.

The broad catalogue target remains unresolved. The restricted models are illustrative equilibrium and transition calculations, not estimates or universal conclusions about microfinance effectiveness.

\section{Completion Estimate}
45\%

This is a post hoc, heuristic estimate for the full catalogue target.
The attempt proves assignment-law dependence under interference, solves an entry-market benchmark, and derives post-credit asset-persistence value and market-policy identification. Full microfinance transformation still requires supported long-horizon saturation responses, tracked closures and movers, and a national transition with endogenous lending, default, funding, wages, and demand.

\begin{thebibliography}{9}
\bibitem{micro} A. Banerjee, E. Duflo, R. Glennerster, and C. Kinnan (2015).
\emph{The Miracle of Microfinance? Evidence from a Randomized Evaluation}.
\url{https://www.aeaweb.org/articles?id=10.1257/app.20130533}.
The primary publication supplies the randomized-access context; no numerical empirical findings are imported into the models above.
\end{thebibliography}

\end{document}
